% !TeX root = ../../../../main.tex
% chapters/sections/chapter5/subsections/a_shock/delta.tex

\paragraph{価格分散\(\Delta_t^H\)の導入による厚生損失についての順位}
\label{sec:chapter5_a_shock_delta}
効用と厚生（\(\Delta_t^H\)無し）の図\ref{fig:chapter5_a_shock_graphs_utility_H_without_delta}と効用と厚生（\(\Delta_t^H\)有り）の図\ref{fig:chapter5_a_shock_graphs_utility_H_with_delta}を比較すると、補正による価格分散\(\Delta_t^H\)の導入が各政策の評価にどのような影響を与えるかがわかる。
この導入による厚生\(W_s^H\)の下落幅が大きい順に政策を示す。
\begin{table}[H]
\centering
\caption{価格分散\(\Delta_t^H\)の導入による厚生損失についての順位}
\label{tab:chapter5_a_shock_delta}
\begin{tabular}{lccc}
\hline
金融政策 & \(W_s^H\)（\(\Delta_t^H\)無し） & \(W_s^H\)（\(\Delta_t^H\)有り） & 下落幅（損失） \\
\hline
実質消費水準目標（CLT） & $-1.3284$ & $-1.3289$ & 0.0005 \\
生産水準目標（OLT） & $-1.3284$ & $-1.3289$ & 0.0005 \\
生産者物価テイラールール（TR-PPI） & $-1.5139$ & $-1.5144$ & 0.0005 \\
名目総消費水準目標（NCLT） & $-1.3399$ & $-1.3403$ & 0.0004 \\
名目GDP水準目標（NGDPLT） & $-1.3399$ & $-1.3403$ & 0.0004 \\
潜在生産水準目標（POLT） & $-1.5476$ & $-1.5480$ & 0.0004 \\
消費者物価水準目標（CPLT） & $-1.3651$ & $-1.3654$ & 0.0003 \\
消費者物価インフレ目標（IT-CPI） & $-1.3654$ & $-1.3657$ & 0.0003 \\
消費者物価テイラールール（TR-CPI） & $-1.5141$ & $-1.5144$ & 0.0003 \\
生産者物価インフレ目標（IT-PPI） & $-1.4925$ & $-1.4925$ & 0.0000 \\
生産者物価水準目標（PPLT） & $-1.5486$ & $-1.5486$ & 0.0000 \\
\hline
\end{tabular}
\end{table}
まずグラフや表から読み取れるように、すべての政策において\(\Delta_t^H\)無しの場合よりも\(\Delta_t^H\)有りの場合の方が厚生\(W_s^H\)が低下している。
これは価格分散の乖離\(\hat{\Delta}_t^H\)が定義上1以上の値をとり、式\eqref{form:chapter3_representative_household_welfare_home_representative_W_H_def}\eqref{form:chapter5_a_shock_welfare_u_H_minus_u_H_ss_eq}により、これが大きいほど厚生\(W_s^H\)が押し下げられることになるからである。